\section{Conclusion [EARLY DRAFT // WORK IN PROGRESS]}
\label{sec:conclusion}

\subsection{Summary and Discussion}
\label{sec:summary}

We have analyzed the network structure of \module{Mathlib} at three levels of granularity plus their cross-level interaction. Table~\ref{tab:summary} collects the key quantitative findings across all four columns---module, declaration, namespace, and cross-level---extending the two-level comparison of Table~\ref{tab:thm-vs-module} with the namespace graph of \S\ref{sec:namespace-graph} and the module--declaration interaction data of \S\ref{sec:module-declaration}.

\begin{table}[ht]
\centering
\caption{Cross-level summary of structural findings.}
\label{tab:summary}
\small
\begin{tabular}{lrrrr}
\toprule
Metric & \S\ref{sec:module-import} Module & \S\ref{sec:theorem-premise} Declaration & \S\ref{sec:namespace-graph} Namespace & \S\ref{sec:module-declaration} Cross-Level \\
\midrule
Nodes                      & $7{,}563$        & $308{,}129$      & $10{,}097$       & --- \\
Edges                      & $23{,}570$       & $8{,}436{,}366$  & $332{,}081$      & --- \\
Cross-namespace edges       & $37.1\%$         & $50.9\%$         & $85.8\%$         & --- \\
Intra-module edges          & ---              & $9.6\%$          & ---              & --- \\
Centrality separation       & partial          & complete         & complete         & --- \\
Community--namespace NMI   & ---              & $0.34$           & $0.25$           & --- \\
Modularity                 & ---              & $0.48$           & $0.28$           & --- \\
Single-node removal impact & $29$ components  & $\le 6$ nodes    & GCC $20.8\%$ at $20\%$ & --- \\
Containment (depth $1$)    & ---              & ---              & $22.2\%$         & --- \\
Theorem/lemma in-degree ratio & ---            & $1.47\times$     & ---              & --- \\
Module cohesion (mean)     & ---              & ---              & ---              & $0.107$ \\
Zero-cohesion modules      & ---              & ---              & ---              & $8.4\%$ \\
$G_{\mathrm{file}}/G_{\mathrm{module}}$ edge ratio & ---  & ---  & ---              & $9.1\times$ \\
Active imports             & ---              & ---              & ---              & $72.1\%$ \\
Indirect declaration deps  & ---              & ---              & ---              & $92.2\%$ \\
\bottomrule
\end{tabular}
\end{table}

The four levels of analysis tell a convergent story from different vantage points.

At the module level (\S\ref{sec:module-import}), the module tree~$T$ and the module graph~$G_{\mathrm{module}}$ exhibit a systematic divergence: $37.1\%$ of reduced import edges cross namespace boundaries, and $17.5\%$ of all import edges are transitively redundant---added by developers for human organizational convenience rather than inherited mathematical necessity (\S\ref{sec:transitive-reduction}). The deep funnel topology ($153$ layers; \S\ref{sec:dag-depth}) and the three-way divergence of centrality measures (\S\ref{sec:centrality}) reveal a library whose architecture is shaped by both logical constraint and human organizational practice. We proposed that \module{Mathlib} has evolved into a \emph{curated small-world}: high local clustering within namespaces, short global paths through deliberate bridge modules, and a maintenance ecology that stabilizes this structure against both accidental fragmentation and unchecked growth (\S\ref{sec:module_conclusion}).

At the namespace level (\S\ref{sec:namespace-graph}), the family of namespace dependency graphs $G_{\mathrm{ns}}^{(k)}$ reveals that the organizing power of human-imposed boundaries erodes rapidly with granularity. Containment---the proportion of edges staying within namespace boundaries---drops from $48.9\%$ at the top-level directory to $22.2\%$ at depth~$1$ and saturates at ${\sim}13\%$ by depth~$3$ (Table~\ref{tab:containment-decay}). The steepest drop occurs at the transition from discipline-level groupings to sub-discipline namespaces, confirming that the broad taxonomic categories of mathematics retain genuine structural force while finer naming distinctions do not. The $14$-percentage-point gap between cross-directory rates at the import level ($37.1\%$) and the declaration level ($51.1\%$) shows that the module graph systematically underestimates the true extent of cross-disciplinary dependency (Remark~\ref{rem:import-vs-decl-cross}).

At the declaration level (\S\ref{sec:theorem-premise}), finer resolution exposes structure that module-level analysis conceals. The hub structure decomposes into two qualitatively distinct layers---a \emph{language infrastructure} layer (\decl{OfNat.ofNat}, \decl{Eq.refl}, \decl{DFunLike.coe}) determined by Lean's type system, and a \emph{mathematical infrastructure} layer (\decl{CategoryTheory.Category}, \decl{Real}, \decl{TopologicalSpace}) determined by the logical structure of mathematics itself (Remark~\ref{rem:two-layers}). Cross-boundary density rises from~$37.1\%$ (cross-directory, at the import level) to~$50.9\%$ (cross-namespace, at the declaration level; \S\ref{sec:thm-cross-namespace}), revealing that module boundaries create an illusion of disciplinary containment that dissolves at finer granularity. Community detection recovers seven recognizable mathematical disciplines from the dependency topology alone (Table~\ref{tab:thm-communities}), with modularity~$0.48$, but their alignment with the namespace hierarchy is only moderate (NMI~$= 0.34$; Table~\ref{tab:thm-nmi}). Among major namespaces, only \ns{CategoryTheory} achieves near-complete community correspondence ($97.4\%$), suggesting genuine disciplinary autonomy. The $44.8\%$ leaf theorem rate (\S\ref{sec:thm-leaves}) marks the boundary between the library's internal reasoning structure and its external API surface, with the \texttt{@[to\_additive]} machinery leaving a distinctive structural fingerprint of systematically generated content.

At the cross-level (\S\ref{sec:module-declaration}), the overlay of all three organizational structures yields the sharpest results. Module cohesion of~$0.107$ (\S\ref{sec:cohesion}) and namespace containment of $22.2\%$ at depth~$1$ (\S\ref{sec:containment-decay}) confirm that both file and namespace boundaries are porous to logical dependencies. The module graph functions as a sparse gateway: $G_{\mathrm{file}}$ has $9.1\times$ more edges than~$G_{\mathrm{module}}$, with $92.2\%$ of cross-file dependencies reached through transitive chains rather than direct imports (\S\ref{sec:import_vs_declaration}). Yet the two human organizational systems---naming and filing---are internally coherent: the NMI between namespaces and modules is~$0.71$, far exceeding the NMI of~$0.34$ between namespaces and dependency-based communities (\S\ref{sec:ns-mod-cross}). The divergence between human organizational process and inherited mathematical structure is thus systematic, not accidental: human decisions are self-consistent (NMI~$0.71$), but the cross-cutting nature of mathematical dependency resists any hierarchical classification (cohesion~$0.107$, containment~$22.2\%$).

\paragraph{The interplay of product and process.}
\label{sec:interplay}
The central claim of this paper is that the network structure of a formal mathematical library records not only the inherited mathematical structure, but also the human organizational process. Our data substantiate this claim with five process-traces embedded in the product, corresponding to the five observations of \S\ref{sec:contribution}:

\begin{enumerate}
\item \textbf{Organizational divergence.} The two human organizational systems---file layout and namespace naming---are internally coherent (NMI~$= 0.71$; \S\ref{sec:ns-mod-cross}) but poorly aligned with dependency-based community structure (NMI~$= 0.34$; \S\ref{sec:sa-community}). Module cohesion of~$0.107$ (\S\ref{sec:cohesion}) and containment decay from $48.9\%$ to ${\sim}13\%$ (\S\ref{sec:containment-decay}) confirm that both file and namespace boundaries are porous to logical dependencies. The hub structure of $G_{\mathrm{thm}}$ reveals a further divergence: the in-degree ranking is dominated by language infrastructure---numeral coercions, equality reasoning---rather than mathematical content (Remark~\ref{rem:two-layers}), recording the design of Lean~4's type system as a distinct layer separable from the mathematical product.

\item \textbf{Cognitive compression.} The module graph compresses the declaration-level dependency network by $9.1\times$ in edges (\S\ref{sec:import_vs_declaration}): $92.2\%$ of cross-file declaration dependencies are reached through transitive import chains rather than direct imports. At the declaration level, the bounded out-degree distribution (lognormal, max~$522$; Table~\ref{tab:thm-powerlaw}) encodes a constraint on the process---no single proof invokes an unbounded number of premises---while the heavy-tailed in-degree ($\alpha \approx 1.78$) reflects the open-ended accumulation of citations in the product. The module graph is a sparse cognitive summary; the declaration graph records a different kind of cognition, shaped by proof decomposition habits and bounded working memory.

\item \textbf{Structural redundancy.} The $17.5\%$ transitive redundancy in $G_{\mathrm{module}}$ (\S\ref{sec:transitive-reduction}) records the gap between what the compiler requires and what the human developer requires. These edges follow the contours of the module tree---developers import along cognitive categories, not along minimal dependency paths. The redundancy is directional: it inflates out-degree without affecting in-degree, confirming that it originates in the importer's cognitive needs.

\item \textbf{Topological inversion.} The three layers exhibit strikingly different topologies (\S\ref{sec:dag-depth}, \S\ref{sec:thm-dag-depth}). The module graph forms a deep funnel ($153$ layers), the declaration graph is shallow and top-heavy ($84$ layers, $38.8\%$ sources at layer~$0$), and the namespace graph is nearly flat ($8$ layers after SCC condensation). Human-organized structure is the deepest; logical structure is the shallowest. This inversion records the fact that the file hierarchy is shaped by incremental, convention-driven layering, while mathematical logic resists deep stratification.

\item \textbf{Semantic flattening.} Formalization preserves traditional labels but flattens their meaning. Declarations marked \texttt{theorem} have $1.47\times$ the mean in-degree of those marked \texttt{lemma} (Table~\ref{tab:thm-vs-lemma}), partially validating human importance judgments in aggregate---but contradicting them in many individual cases (\S\ref{sec:thm-vs-lemma}). Meanwhile, the $44.8\%$ zero-citation rate among theorems (\S\ref{sec:thm-leaves}) includes a substantial fraction of \texttt{@[to\_additive]} mirrors that exist for API completeness rather than logical necessity (Remark~\ref{rem:leaf-theorems}). These community norms inflate the graph's surface area without deepening its logical content.
\end{enumerate}

But the relationship between product and process is not unidirectional. The product reshapes the process in return, through at least two mechanisms. First, the $37.1\%$ cross-directory edge rate in $G_{\mathrm{module}}^{-}$ (\S\ref{sec:cross-namespace}) exerts reorganization pressure on the directory hierarchy: when modules in one namespace systematically depend on another, developers eventually redraw the cognitive categories---as witnessed by \module{Mathlib}'s ongoing absorption of \module{RingTheory} into \module{Algebra}. Second, the deep dependency chain ($153$ layers; \S\ref{sec:dag-depth}) constrains the development process directly: modifications to high-PageRank modules trigger recompilation cascades that shape contributor behavior, incentivizing stability at the core and experimentation at the periphery. Formalization thus initiates a feedback cycle: the module tree shapes the module graph; the module graph reshapes the module tree. Network analysis, as developed in this paper, provides the quantitative vocabulary for tracing this cycle.

\subsection{Implications for Practice}
\label{sec:practice}

The metrics developed in this paper are not merely descriptive; several admit direct operational use.

\paragraph{Library maintenance.} Module cohesion (Definition~\ref{def:cohesion}) can be computed incrementally as new declarations are added. A sustained decline in cohesion for a module signals that its scope has drifted from its original organizational intent. The zero-citation rate (\S\ref{sec:thm-leaves}) identifies leaf declarations that may represent redundant API surface; periodic audits could reduce the maintenance burden on the approximately $2{,}000$ pending pull requests that \module{Mathlib} currently faces.

\paragraph{Compilation and CI.} The PageRank and betweenness rankings identify modules whose modification triggers the widest recompilation cascades. Integrating these rankings into continuous integration pipelines---for example, by assigning higher test priority to pull requests that touch high-centrality modules---could reduce the latency between contribution and merge.

\paragraph{Language design.} The two-layer hub structure (\S\ref{sec:thm-types}) suggests that future proof assistants could benefit from an explicit architectural boundary between language infrastructure and mathematical content, allowing dependency analysis and compilation to target each layer independently. The $92.2\%$ indirect dependency rate (\S\ref{sec:import_vs_declaration}) further suggests that finer-grained import mechanisms---at the declaration level rather than the module level---could reduce both redundancy and recompilation scope, though Lean's typeclass inference mechanism, which relies on broad import scopes, would need to be accommodated.

\paragraph{AI training data.} For systems that learn premise selection or proof generation from \module{Mathlib}, our analysis reveals structural features that can serve as training metadata: the theorem/lemma annotation provides a human importance signal; the Louvain community label provides a disciplinary tag; and the declaration type (theorem vs.\ definition vs.\ axiom) distinguishes the roles that different nodes play in the dependency network. The two-layer hub structure further suggests that language-infrastructure nodes may need to be treated differently from mathematical-infrastructure nodes in a learning pipeline. Whether topological depth in $G_{\mathrm{thm}}$---the length of the longest dependency chain terminating at a given declaration---correlates with proof difficulty remains an open question: depth measures logical dependency chains, not the cognitive or computational effort required to construct a proof. If such a correlation holds, proof generation systems might benefit from curriculum strategies that ascend the chain incrementally, but this hypothesis requires empirical validation.

\paragraph{Monitoring structural drift.} As AI-generated proofs become a larger fraction of \module{Mathlib} contributions, the metrics established here---redundancy rate, cohesion distribution, degree distribution exponents, cross-namespace density---provide a dashboard for detecting structural drift. A sudden change in any of these indicators would signal that the library's character is shifting in ways that merit human review.

\subsection{Limitations and Future Work}
\label{sec:limitations}

Several limitations constrain the scope of our conclusions.

\emph{Static snapshot.} All statistics are computed from a single commit (\texttt{534cf0b}, February 2026). We cannot determine whether the structural properties we observe---the $17.5\%$ redundancy rate, the $153$-layer depth, the $0.107$ cohesion---are stable, growing, or approaching critical thresholds. The evolutionary hypotheses advanced in \S\ref{sec:module_conclusion} remain untested predictions.

\emph{Social process unobserved.} The ``process'' side of our analysis is inferred from structural traces in the product, not observed directly. We do not analyze PR review workflows, contributor patterns, commit history, or the social dynamics of the \module{Mathlib} community. The curation mechanisms that maintain the library's small-world topology (\S\ref{sec:module_conclusion}) are posited on structural evidence but not empirically verified through process data.

\emph{Correlation, not causation.} We measure the divergence between the module tree~$T$ and the dependency graphs~$G_{\mathrm{module}}$ and~$G_{\mathrm{thm}}$, but we cannot disentangle the causal direction. Does~$T$ shape~$G$ (developers import along cognitive categories), or does~$G$ shape~$T$ (logical dependencies force namespace reorganization)? Our data are consistent with bidirectional influence (\S\ref{sec:interplay}), but a causal account would require longitudinal or experimental evidence that this study does not provide.

\paragraph{Future directions.}
\label{sec:future}

\emph{Longitudinal evolution.} Repeating this analysis across \module{Mathlib} versions---particularly across the Lean~3 to Lean~4 migration---would test whether structural indicators (redundancy rate, DAG depth, cohesion, cross-namespace density) are stable invariants of mathematical content or contingent features of a particular proof assistant and community. If the $17.5\%$ redundancy rate and $0.107$ cohesion are increasing over time, this would provide evidence for the critical-depth hypothesis of \S\ref{sec:module_conclusion}; if they are stable, it would suggest that the library has reached a structural equilibrium.

\emph{AI-induced structural drift.} As AI systems contribute an increasing fraction of \module{Mathlib}'s proofs---through tactic automation, premise selection, or end-to-end proof generation---the process-traces identified in \S\ref{sec:interplay} may shift systematically. AI-generated proofs are not subject to the same cognitive complexity bound that shapes the out-degree distribution (item~2 above); they may not follow the namespace-aligned import patterns that produce the $17.5\%$ redundancy (item~3). The current data provide a human-authored baseline against which AI-induced structural drift can be measured. More fundamentally, an AI-dominant production process could alter the granularity equilibrium itself (\S\ref{sec:three-layers}): the approximately $44$-declaration-per-module ratio reflects human cognitive capacity, not mathematical necessity. AI contributors, unconstrained by working-memory limits, might generate libraries with radically different module sizes---finer for compilation parallelism or coarser for logical coherence---producing dependency topologies whose statistical profiles differ not because the mathematics differs, but because the producer differs. The structural snapshot we report is, in this sense, a snapshot of a particular production regime in formal mathematics.

\emph{Deeper namespace analysis.} The containment decay curve of \S\ref{sec:containment-decay} characterizes how namespace boundaries lose organizing power as granularity increases. Natural extensions include community detection on the family $G_{\mathrm{ns}}^{(k)}$ at each depth---testing whether dependency-based communities align more closely with the namespace hierarchy at coarser depths---and centrality analysis of the namespace graph, which would identify which namespaces serve as structural bridges between mathematical disciplines.

\emph{Premise co-occurrence analysis.} After filtering high-frequency infrastructure pairs, the co-occurrence graph may reveal clusters of premises that are cognitively associated---``proof toolkits'' that practitioners reach for together, independent of logical necessity. Such clusters would constitute the most direct trace of proof-search heuristics in the product, complementing the structural traces analyzed in this paper. Comparing these clusters with the communities detected in $G_{\mathrm{thm}}$ (\S\ref{sec:thm-community}) would test whether co-usage patterns align with dependency-based communities or reveal an orthogonal organizational dimension.

\emph{Cross-library comparison.} Applying the same analysis to other formal libraries---Isabelle/AFP, Coq/MathComp, Mizar/MML---would distinguish structural features intrinsic to mathematics from those contingent on a particular proof assistant or community culture. If the funnel topology, the two-layer hub structure, and the near-zero module cohesion appear across libraries built on different type theories and maintained by different communities, they are properties of formalized mathematics itself; if they diverge, they are artifacts of the medium.

\emph{Product or process?} A deeper question underlies all of these directions: is the dependency topology we observe a property of \emph{mathematics}, or of \emph{how mathematics is currently produced}? If AI systems, operating under different cognitive constraints, were to build an equivalent library from scratch, would the resulting graph share the same degree distributions and community structure? The metrics developed in this paper---redundancy rate, cohesion, containment, centrality rankings, depth profiles---provide the vocabulary in which the answer, when it comes, can be expressed.
